\documentclass{report}
\usepackage{amsmath}
\usepackage{amssymb}
\usepackage{hyperref}
\usepackage[margin=1in]{geometry}
\setlength{\parindent}{0pt}

\begin{document}
\thispagestyle{empty}
\begin{center} \LARGE Assignment 8 \end{center}

\section*{Part 1: Mathematical Derivations}

\subsection*{Q1}
\begin{enumerate}
    \item we should notice that the standard deviation for \(q(x_t\vert x_{t-1} ) = \sqrt{\beta_t}\). the mean is also \(\sqrt{1-\beta_t}x_{t-1}\) and so \(x_t = \sqrt{1 - \beta_t} x_{t-1} + \sqrt{\beta_t} \epsilon_t\) where \( \epsilon_t \sim \mathcal{N}(0, I)\)
    \item given \(\alpha_t\), we can rewrite as 
    \begin{equation}
        \begin{split}
            x_t &= \sqrt{\alpha_t} x_{t-1} + \sqrt{1 - \alpha_t} \epsilon_t\\
            x_{t-1} &= \sqrt{\alpha_{t-1}} x_{t-2} + \sqrt{1 - \alpha_{t-1}} \epsilon_{t-1}
        \end{split}
    \end{equation}
    which gives 
    \begin{equation}
        x_t = \sqrt{\alpha_t \alpha_{t-1}} x_{t-2} + \sqrt{\alpha_t (1 - \alpha_{t-1})} \epsilon_{t-1} + \sqrt{1 - \alpha_t} \epsilon_t
    \end{equation}
    since variances sum, we get 
    \begin{equation}
        x_t = \sqrt{\alpha_t \alpha_{t-1}} x_{t-2} + \sqrt{1 - \alpha_t \alpha_{t-1}} \epsilon_{t-2}'
    \end{equation}

    unfolding all the way to \(x_0\), we get

    \begin{equation}
        x_t = \sqrt{\bar{\alpha}_t} x_0 + \sqrt{1 - \bar{\alpha}_t} \epsilon, \quad \text{where } \epsilon \sim \mathcal{N}(0, I)
    \end{equation}

    which is the same as 
    $$\mathbf{q(x_t | x_0) = \mathcal{N}\left(x_t; \sqrt{\bar{\alpha}_t} x_0, (1 - \bar{\alpha}_t) I\right)}$$

    \item the distribution just approaches the standard gaussian distribution \(\mathcal N(0,I)\). as \(\bar\alpha_T\) goes to zero the mean goes to zero and the variance goes to \(I\)

\end{enumerate}

\subsection*{Q2}

\begin{enumerate}
    \item \begin{equation}
        \log p_\theta(x) = \log \left( \frac{e^{-f_\theta(x)}}{Z_\theta} \right) = -f_\theta(x) - \log Z_\theta
    \end{equation}

    \begin{equation}
        s_\theta(x) = \nabla_x \log p_\theta(x) = \nabla_x \left( -f_\theta(x) - \log Z_\theta \right)
    \end{equation}
    but note that \(Z_\theta\) is constant wrt to \(x\) so the derivative is zero and thus
    \begin{equation}
        s_\theta(x) = -\nabla f_\theta(x)
    \end{equation}

    but note that \(Z_\theta\) is constant wrt to \(x\) so the derivative is zero.
    \item solving for epsilon from part 1,
    \begin{equation}
        \epsilon = \frac{x_t - \sqrt{\bar{\alpha}_t}x_0}{\sqrt{1 - \bar{\alpha}_t}}
    \end{equation}
    substituting
    \begin{equation}
        q(x_t \mid x_0) = \frac{1}{(2\pi)^{d/2} |(1 - \bar{\alpha}_t)\mathbf{I}|^{1/2}} \exp\left( -\frac{\|x_t - \sqrt{\bar{\alpha}_t}x_0\|^2}{2(1 - \bar{\alpha}_t)} \right)
    \end{equation}


    and taking the log gives
    \begin{equation}
        \log q(x_t \mid x_0) = -\frac{\|x_t - \sqrt{\bar{\alpha}_t}x_0\|^2}{2(1 - \bar{\alpha}_t)} + C
    \end{equation}

    grad of this gives 
    \begin{equation}
        \begin{split}
            \nabla_{x_t} \log q(x_t \mid x_0) &= -\frac{2(x_t - \sqrt{\bar{\alpha}_t}x_0)}{2(1 - \bar{\alpha}_t)}\\
            &= -\frac{x_t - \sqrt{\bar{\alpha}_t}x_0}{1 - \bar{\alpha}_t}\\
            &= -\frac{\sqrt{1 - \bar{\alpha}_t}\epsilon}{1 - \bar{\alpha}_t} \\ 
            & = -\frac{\epsilon}{\sqrt{1 - \bar{\alpha}_t}}
        \end{split}
    \end{equation}
    isolating for \(\varepsilon\),
    \begin{equation}
        \epsilon = -\sqrt{1 - \bar{\alpha}_t} \nabla_{x_t} \log q(x_t \mid x_0)
    \end{equation}
    \item to know the exact probability would require calculating \(Z\) exactly which is hard. 
\end{enumerate}

\subsection*{Q3}

\begin{enumerate}
    \item \begin{enumerate}
        \item \(\varepsilon\) is gaussian noise, 
        
        \(t\) is uniform in \([t]\).
        \item the network input is a noisy image \(x_t\) and a timestep \(t\). 
        \item the network is trying to learn the gaussian noise vector that was added to \(x_0\) to get \(x_t\). MSE penalizes the difference between the prediction and the true noise.
    \end{enumerate}
    \item \begin{enumerate}
        \item \(\sigma_t\) controls the variance of the component added at each reverse step. 
        \item at \(t=1\), this is the final prediction. this step is deterministic of hte input from the prior step
    \end{enumerate}
    \item 
\end{enumerate}



\section*{Reflection + Statistics Section}
(a)

(b) 

(c) 

(d) 
Contributions: \\

Total: \\

Part 1: \\
Q1: \\

Part 2: \\
Q1: \\

Part 3: \\
Q1: \\

(e)


\end{document}
